\section{Ridge Regression}
\label{sec:ridge}

\subsection{Ordinary least squares}

Given a design matrix $X \in \mathbb{R}^{n\times p}$ (rows are
(country, year) observations, columns are features) and target vector
$y \in \mathbb{R}^n$, ordinary least squares (OLS) solves

\begin{equation}
\label{eq:ols}
\hat\beta_{\mathrm{OLS}} = \arg\min_{\beta} \left\| y - X\beta \right\|_2^2
\end{equation}

with closed-form solution
$\hat\beta_{\mathrm{OLS}} = (X^\top X)^{-1}X^\top y$, provided
$X^\top X$ is invertible.

\subsection{The problem with correlated features}

This project's feature set includes multiple lags and rolling statistics
of the same underlying series (\texttt{co2\_lag1}, \texttt{co2\_lag3},
\texttt{co2\_lag5}, \texttt{co2\_lag10}, \texttt{co2\_roll5\_mean},
\texttt{co2\_roll10\_mean}), which are, by construction, highly correlated
with one another. When columns of $X$ are strongly correlated,
$X^\top X$ is close to singular, and $\hat\beta_{\mathrm{OLS}}$ becomes
highly sensitive to small changes in the data: its variance inflates even
though it remains unbiased.

\subsection{Ridge regression}

Ridge regression adds an $\ell_2$ penalty on the coefficients:

\begin{equation}
\label{eq:ridge}
\hat\beta_{\mathrm{ridge}} = \arg\min_{\beta}
\left\| y - X\beta \right\|_2^2 + \alpha \left\| \beta \right\|_2^2,
\qquad \alpha \ge 0
\end{equation}

with closed-form solution

\begin{equation}
\label{eq:ridge-closed-form}
\hat\beta_{\mathrm{ridge}} = \left(X^\top X + \alpha I\right)^{-1} X^\top y
\end{equation}

Adding $\alpha I$ to $X^\top X$ in Equation~\ref{eq:ridge-closed-form}
guarantees invertibility for any $\alpha>0$ and shrinks the coefficients of
correlated features toward one another and toward zero. This trades a
small amount of bias for a reduction in variance \cite{hoerl1970ridge},
and is the reason this project prefers Ridge over OLS given its
correlated lag and rolling-statistic feature set.

\subsection{Preprocessing}

Because the ridge penalty in Equation~\ref{eq:ridge} is applied uniformly
to all coefficients, features must be on comparable scales before
fitting; unlike tree-based models (Section~\ref{sec:ensembles}), Ridge
requires feature scaling, and that scaler must be fit on training rows
only, to avoid leaking validation or test distributional information into
its parameters.
